\chapter{Related Work}
\label{ch:related}

The literature on honeyword-based authentication has two threads. The
first is the design of \emph{generators} that produce honeywords
intended to be difficult to distinguish from real passwords. The
second, more recent, is the design of \emph{distinguishers} that
attempt to break flatness under stronger attacker models than the ones
generators are usually evaluated against. Both threads inform this
thesis: the ten generators of \Cref{sec:method-gens} are drawn from
the first thread, and the CNN of \Cref{sec:method-cnn} is drawn from
the second. This chapter reviews both, and \Cref{sec:related-position}
describes how the filtering construction of \Cref{sec:method-filter}
sits between them.

\section{Heuristic Honeyword Generators}
\label{sec:related-heuristic}

The earliest honeyword generators are the four \emph{chaffing}
constructions proposed alongside the primitive by
\textcite{juels2013honeywords}: \emph{chaffing-by-tweaking},
\emph{chaffing-by-tail-tweaking}, \emph{chaffing-with-a-password-model},
and \emph{chaffing-with-tough-nuts}. Chaffing-by-tweaking applies
character-class substitutions to the real password to produce visually
similar decoys; chaffing-by-tail-tweaking restricts the edits to a
short suffix. Chaffing-with-a-password-model draws candidates from an
external password model such as a PCFG or a Markov chain, so that the
decoys look like plausible passwords in the aggregate rather than like
edits of the real one. Chaffing-with-tough-nuts inserts random,
high-entropy strings that no attacker would submit at login and that
therefore serve as guaranteed alarms.

These constructions were evaluated against attackers that scored
candidates independently, typically PCFG-based rankers of the kind
introduced by \textcite{weir2009pcfg} and Markov rankers of the kind
introduced by \textcite{narayanan2005fast}. The flatness numbers
reported under this evaluation model are close to $1/k$ for
carefully-tuned chaffing schemes, which motivated the follow-up
literature. Two of these constructions, chaffing-by-tweaking and
chaffing-with-a-password-model, appear in our benchmark
(\Cref{sec:method-gens}) as reference points against the more
sophisticated generators discussed below, together with a hybrid
chaffing method that mixes the two and that we introduce for
comparison.

The main line of criticism against the heuristic generators is
\emph{per-password plausibility}: the decoys are designed so that each
one looks like a plausible password in isolation, but there is no
constraint on how they look \emph{as a set}. A generator can produce
$k-1$ plausible but distributionally-clustered decoys that leave the
real password stranded, and the attacker only needs to notice the
cluster to identify the outlier. This is the phenomenon that later
generators, and the filter of \Cref{sec:method-filter}, are designed
to address.

\section{Representation-Learning Approaches: HoneyGen}
\label{sec:related-honeygen}

\textcite{dionysiou2021honeygen} propose \emph{HoneyGen}, a generator
that combines a FastText password model with rule-based tweaking. The
password model produces nearest-neighbor candidates for the real
password in the FastText embedding space (\Cref{sec:bg-fasttext}), and
a tweaking module produces additional candidates by character-class
substitution. The two streams are merged, pruned using a
frequency-based filter that removes overly common candidates, and the
top $k-1$ are retained as honeywords.

The evaluation in the original paper is against frequency-based and
PCFG-based attackers, and HoneyGen achieves flatness scores that
improve on the chaffing baselines. Under this class of attacker the
paper's conclusion is uncontroversial: representation learning
captures more of the password distribution than either rules or PCFGs
alone, and the resulting decoys are harder to distinguish under a
ranker of the kind those attackers rely on.

Our position with respect to HoneyGen is precise. We reuse the
generator, without modification, as the source of candidates for six
of the ten methods in the benchmark (HoneyGen Baseline, Set-Filter,
Prescreen-Add, Prescreen-Ratio, Hybrid-Add, Hybrid-Ratio). What we do
not reuse is the paper's evaluation: our attacker is the CNN of
\Cref{sec:bg-cnn} rather than a frequency ranker, and the flatness
scores we report under that attacker are substantially worse than the
ones the original paper reports (\Cref{ch:results}). The four filtered
variants are the mechanism by which we recover flatness against this
stronger attacker without modifying the generator itself.

\section{Large-Language-Model Approaches}
\label{sec:related-llm}

More recent work uses large language models as the source of
honeyword candidates. The chunk-level construction of
\textcite{yu2023honeychunk} has drawn particular attention: a password
is disassembled into semantic chunks that a pre-trained language model
perturbs one at a time, and candidate sets are drawn with
Metropolis--Hastings sampling -- an accept-or-reject sampling scheme
that steers the generated candidates toward regions the model assigns a
high real-password probability -- while locality-sensitive hashing, a
technique that maps similar strings to the same bucket, is used to drop
near-duplicates from the pool. These are among the most linguistically
sophisticated generators the honeyword literature has produced, and
they report clear gains in semantic plausibility over both HoneyGen and
GAN-based methods.

The evaluation in this line, however, follows the same statistical
plausibility paradigm as the earlier work. Flatness is measured
against ranker-style attackers, and the metrics reported are
distributional distances rather than the success rate of a neural
distinguisher trained specifically to identify the real password. The
gap this leaves is exactly the one \textcite{dani2024passfilter}
demonstrated for HoneyGen: a generator can be indistinguishable under
one attacker class and yet break under a stronger one. We do not
evaluate an LLM-based generator directly in this thesis, since our
focus is on the reuse of the neural distinguisher as a defensive
oracle rather than on the generation side. The construction of
\Cref{sec:method-filter} is agnostic to the source of candidates and
would apply to an LLM-based generator without modification;
\Cref{ch:conclusion} returns to this point.

\section{GAN-Based Honeyword Generators}
\label{sec:related-gan}

\textcite{hitaj2019passgan} adapt the generative adversarial network
of \textcite{goodfellow2014generative} to password modeling, in the
attack setting: the goal is to produce guess sequences that recover
more of a leaked corpus than PCFG or Markov models produce at the same
budget. The adaptation is to train the generator against a
discriminator that tries to separate generated strings from real
leaked ones, with both networks operating on character sequences.

GANs have already been studied as a source of honeyword candidates.
\textcite{fauzi2020passgan} analyse a game in which the
defender uses the generator of a PassGAN model to produce honeywords
while the attacker uses its discriminator to pick out the real
password, and assess the feasibility of PassGAN as a honeyword
generation strategy. We include a GAN in this same spirit as one of the
generators in our benchmark: the intuition is that a model that has
learned the password distribution should, in principle, produce decoys
that look like real passwords. Our GAN is a character-level generator
trained from scratch on the same corpus as the other generators (a
\emph{self-trained} GAN), and we sample decoys from it directly.
GANs are prone to \emph{mode collapse}, where the
generator learns to emit only a narrow set of outputs instead of
covering the full variety of the training distribution; for passwords
this shows up as the same few common strings being produced over and
over. To obtain enough distinct decoys despite this, we oversample the
generator -- drawing roughly five times as many candidates as a set
needs -- and deduplicate the resulting pool before assembling each
sweetword set. This GAN is one of the ten generators evaluated in our
benchmark (\Cref{ch:method}).

The theoretical case for using a GAN as a honeyword generator is
subtle, and our results show it does not hold in practice. A generator
that had truly converged to the real password distribution would
produce indistinguishable decoys almost by definition, but real
password distributions have long low-probability tails that GANs are
known to underrepresent, so the resulting decoys cluster in the
high-probability region and the real password, often itself a
low-probability string, is left looking anomalous. In
\Cref{ch:results} we show exactly this: under the CNN attacker the
GAN's real password is pushed to the bottom of the ranking
($F_{\text{rank}} \approx 0.92$), so its apparently low attack success
is an artifact of that inversion rather than genuine flatness, and a
mixed-corpus attacker exposes it. This is why the GAN's flatness
numbers must be read with care rather than taken at face value.

\section{Neural Distinguishers as Attackers: PassFilter}
\label{sec:related-passfilter}

The line of work most directly relevant to this thesis is PassFilter
\parencite{dani2024passfilter}. It trains a character-level
convolutional neural network as a binary classifier that predicts,
given a candidate password, whether it is real or a honeyword. The
network is small by modern standards (two convolutional layers over
character embeddings, followed by a fully connected layer), and it is
trained on labeled pairs of real passwords and honeywords produced by
existing generators.

The main result of the paper is that this classifier violates flatness
on every honeyword scheme they test, including HoneyGen and its
variants, on multiple leaked corpora. The attacker success rate is
substantially higher than $1/k$ across the board, sometimes by an
order of magnitude, and the ranking induced by the classifier
concentrates on the real password with a probability the classical
flatness definition explicitly forbids.

Using a learned discriminator as the attacker is not unique to
PassFilter. As noted in \Cref{sec:related-gan},
\textcite{fauzi2020passgan} already turn the discriminator of
a PassGAN model against its own generator's honeywords; PassFilter can
be read as the general form of the same idea, a classifier trained to
separate real passwords from honeywords regardless of which generator
produced them. Both establish the neural distinguisher as the attacker
class that honeyword schemes must withstand. This thesis takes that
distinguisher and moves it to the defensive side.

Our reading of the PassFilter result is that the honeyword literature
has been evaluating flatness against too weak an attacker class. The
constructive response we develop is to admit the same class of
classifier into the threat model, formalize flatness against it
(\Cref{sec:method-formalism}), and then reuse the classifier itself,
without architectural change, as the enforcement mechanism during
sweetword construction (\Cref{sec:method-filter}). The name we give to
this reuse, \emph{defensive oracle}, is meant to underline that the
same object is used offensively in PassFilter and defensively in this
thesis; only the surrounding pipeline changes.

Two aspects of the PassFilter setup differ from ours in ways that
matter. First, the paper releases model weights that are trained on a
specific corpus and are used to attack sweetwords produced by
generators trained on the same corpus. We do not reuse those weights;
we train our own CNN on a disjoint attacker corpus
(\Cref{sec:method-corpora}), which raises the baseline attacker
success rate but removes any shared-parameter bias between attacker
and generator. Second, the paper evaluates flatness at the level of
individual sets. We report the full ASR curve $\mathrm{ASR}(t)$ over the
guess budget $t$, computed over the accepted sets, for every generator
under a common attacker.

\section{The HBAT Survey}
\label{sec:related-survey}

\textcite{chakraborty2022hbat} survey the honeyword literature and
propose a framework, \emph{honeyword-based authentication techniques}
(HBAT), that organizes the existing schemes along three axes:
generation method, distributed security model, and evaluation
metric. Their treatment of flatness is compatible with, and indeed
formalizes, the definition given by \textcite{juels2013honeywords};
they extend it to $\varepsilon$-flatness and connect it to two further
security notions (DoS resistance and MSV resistance) that we mentioned
in \Cref{sec:bg-flatness}.

The survey's attacker taxonomy runs from frequency-based rankers
through PCFG and Markov models to neural language models. Neural
\emph{distinguishers} of the kind introduced by
\textcite{dani2024passfilter} appear implicitly in the survey's
discussion of adversarial machine learning but are not integrated into
the flatness definition itself. The refinement of the flatness
definition in \Cref{sec:method-formalism} can be read as such an
integration.

\section{Position of This Thesis}
\label{sec:related-position}

The construction developed in this thesis sits at the intersection of
the two threads reviewed above. From the generator thread, it inherits
the ten methods of \Cref{sec:method-gens} without modification. From
the distinguisher thread, it inherits the CNN classifier of
\textcite{dani2024passfilter}, again without architectural change.
What is new is the pipeline in which the two are combined: the
classifier is trained on a disjoint corpus and reused as the oracle
that decides which sweetword sets are stored, and the flatness
conditions the oracle enforces are stated at the level of the set-wide
induced distribution defined in \Cref{sec:method-formalism} rather
than at the level of individual candidates.

Three summary observations situate the contribution.
%
\begin{enumerate}
    \item Prior generators are evaluated against attackers that they
    themselves inspired. Frequency-based generators are evaluated
    against frequency rankers, and representation-learning generators
    are evaluated against Markov or PCFG rankers. This thesis holds the
    attacker class fixed, at the CNN of
    \textcite{dani2024passfilter}, and lets the generators vary; that
    common attacker is what makes the ten methods comparable on a single
    ASR axis. Robustness to the choice of attacker is then examined
    separately (\Cref{sec:results-freq}), by re-evaluating the accepted
    sets against a mixed-corpus CNN and a frequency ranker.

    \item Prior work on neural distinguishers stops at the
    demonstration that flatness is violated. It does not propose a
    construction that restores flatness. This thesis proposes one and
    measures how far it goes.

    \item Prior filtering ideas exist in the honeyword literature, but
    they operate at the candidate level and use non-learned filters
    (frequency pruning, dictionary matching). The two-stage filter of
    \Cref{sec:method-filter} uses a learned filter and, in its
    set-level stage, constrains the distribution rather than the
    individual candidates. The combination is what allows the filter
    to enforce $\varepsilon$-flatness in the sense of
    \Cref{eq:eps-flat}.
\end{enumerate}
%
The remainder of the thesis develops these three points in detail.
\Cref{ch:method} specifies the construction, and \Cref{ch:results}
measures its effect on ASR across the ten generators.
